\section{Discussion}
\label{sec:discussion}

The empirical results in Section~\ref{sec:results} establish two findings that
require interpretation. Direct quantile regression on the HAR structure
delivers a statistically significant upper-tail advantage at short horizons
but is overtaken by simpler quantile constructions --- the HAR-OLS
pseudo-quantile at six months and the GARCH(1,1) Gaussian quantile at twelve
months --- as the horizon extends (RQ1), and the Ni\~no~3.4 climate index
provides no robust forecast improvement at any horizon, its modest
twelve-month full-sample gain being concentrated entirely in the 2022--2024
supply crisis (RQ2). This section
explains why these patterns arise, how they relate to the existing literature
on volatility forecasting and ENSO--commodity linkages, and what they imply for
the calibration of distributional forecasts, the procurement perspective that
motivated the choice of cocoa, and the directions for further research.

%------------------------------------------------------------------------------
\subsection{Why the HAR Family Dominates}
\label{sec:disc-rq1}
%------------------------------------------------------------------------------

The dominance of the HAR family in mean forecasting is structural rather
than incidental. \textcite{corsiSimpleApproximateLongMemory2008}
shows that standard short-memory volatility models, including GARCH, fail to
reproduce the scaling behaviour of volatility under temporal aggregation: their
exponential memory decay erases shocks too rapidly, and multi-step forecasts
converge toward the unconditional variance within a few weeks. The HAR
structure approximates the hyperbolic decay that empirical volatility
exhibits by stacking heterogeneous lags at daily, weekly, and monthly
frequencies, without imposing a fractional integration parameter. This
matters with particular force for cocoa, which
\textcite{alfeusForecastingVolatilityCommodity2022} identify as a long-memory
commodity in which monthly volatility carries the most predictive weight.
Section~\ref{sec:res-mean} reports a correlation of 0.66 between the monthly
HAR component and the twelve-month volatility target, against 0.37 for the
daily component, which is the direct empirical counterpart of Corsi's scaling
argument and explains why HAR-OLS retains its lead even at twelve months while
GARCH(1,1) flatlines.

The upper-tail ranking at long horizons inverts this picture, and the
inversion is informative about what each model actually estimates. The
GARCH quantile forecast (Equation~\ref{eq:garch-quantile}) anchors the
$\tau = 0.95$ boundary to the dispersion of $h$-day average volatility
over the full estimation window, while its conditional mean is updated
daily. At twelve months this construction is conservative but stable: it
does not attempt to extrapolate tail dynamics from the recent past alone.
The rolling HAR-QR quantile, by contrast, is fitted to the empirical
upper tail of a 2,520-day window; when the 2022--2024 crisis pushed
realized volatility far outside the support of that window, the fitted
quantile surface had no comparable episodes to learn from, and the
forecast boundary adjusted only with substantial delay
(Figure~\ref{fig:upper-tail}, panel b2). The result is GARCH's 29\%
quantile-loss advantage over HAR-QR at the twelve-month upper tail
(Table~\ref{tab:upper-qlike-comparison}): a deliberately wide,
slowly-moving quantile loses little in calm decades and avoids the
catastrophic under-prediction that dominates the loss of the more
adaptive models during the crisis. The advantage is one of point loss
only; at twelve months GARCH, HAR-OLS, HAR-QR and HAR-X-QR all belong to
the model confidence set (Section~\ref{sec:res-mcs}), so the ranking
reflects the lowest average loss rather than a statistically separable
edge.

Within the HAR family, the comparison between HAR-QR and the HAR-OLS
pseudo-quantile addresses a more subtle question: does the conditional upper
quantile diverge sufficiently from a shift of the conditional mean to justify
the additional estimation cost of quantile regression?
\textcite{haugomParsimoniousQuantileRegression2016} argue that the gain comes
from two sources: direct quantile estimation is robust to distributional
misspecification, and the monthly volatility component exerts its strongest
influence precisely at the conditional tails. The cocoa results match this
mechanism. The HAR-QR advantage over the HAR-OLS pseudo-quantile is largest at
the short to medium horizons (9\,\% at one week, 7\,\% at one month, 6\,\% at
one day), where volatility dynamics are most state-dependent and the
conditional quantile is furthest from the unconditional residual quantile.
From six months onward the advantage reverses in point terms: HAR-OLS
undercuts HAR-QR by 14\,\% at six months and 24\,\% at twelve months
($0.000786$ vs.\ $0.001041$), though the two remain statistically
indistinguishable in the MCS (Section~\ref{sec:res-mcs}). \textcite{lyocsaImprovingStockMarket2021} report
the same horizon dependence on stock-market data: the value of explicit
quantile estimation is concentrated at short horizons where extreme regime
states drive the conditional distribution. Coverage at the one-day and
one-week horizons is close to the nominal $5\,\%$ ($5.17\,\%$ and $5.39\,\%$).
Christoffersen's $LR_{uc}$ test confirms that unconditional coverage is
indistinguishable from the nominal $5\,\%$ at one day ($p = 0.495$), consistent
with Haugom's finding for equity markets, now extended to an agricultural
commodity. However, the independence component $LR_{ind}$ is rejected at one
day ($p = 0.008$), indicating that violations cluster even at the shortest
horizon: the conditional coverage test is formally rejected ($p = 0.024$).
At longer horizons, overlapping $h$-day targets make consecutive violation
indicators dependent by construction, so the Christoffersen statistics are
computed on the non-overlapping subsequence of every $h$-th forecast
(Table~\ref{tab:christoffersen}, Appendix~\ref{app:supplementary}). On that
basis, unconditional coverage is rejected for HAR-QR at six months, while at
twelve months the non-overlapping subsequence (28 observations) is too short
for an informative test and the raw violation rate of $14.63\,\%$ is the more
telling summary.

%------------------------------------------------------------------------------
\subsection{Interpreting the ENSO Signal}
\label{sec:disc-rq2}
%------------------------------------------------------------------------------

The ENSO result requires less interpretive caution than the research design
anticipated, because the corrected evidence points uniformly in one
direction. Over the full evaluation sample, HAR-X-QR reduces upper-tail
quantile loss by $3.3\,\%$ at the twelve-month horizon --- the only
quantile--horizon combination with a gain of any substance --- while
worsening the forecast at fifteen of the eighteen combinations examined,
and it holds the top MCS rank at no horizon. The robustness diagnostics
in Section~\ref{sec:res-robustness} show that even the twelve-month gain
is fragile: the cumulative gain is negative over the entire pre-crisis
period, materialises wholly within the 2022--2024 supply-crisis window,
follows a lag structure opposite to what an agronomic transmission chain
would predict, and does not replicate on Arabica coffee or raw sugar. The
interpretive task is not to explain away a striking effect, but to
understand why a mechanism with a plausible physical basis leaves so
faint a statistical trace.

The physical channel that motivated the inclusion of ENSO in the first place
is plausibly real but is unlikely to be visible at the lag structure assumed
by the standard climate-yield literature. \textcite{schrothVulnerabilityClimateChange2016}
document that severe El~Ni\~no events in the 1980s coincided with drought
years that materially affected West African cocoa yields, which suggests a
transmission window of six to twelve months from sea-surface anomaly to
harvest impact. Our lag-decay diagnostic shows the opposite pattern:
predictive power peaks at lag~1 and decays sharply by lag~6.
\textcite{berganderDeepLearningCocoa2025} provide the missing link: futures
prices in their analysis reflect roughly $85\,\%$ of the eventual price
impact of a drought shock \emph{before} the actual crop loss occurs. If
professional traders anticipate ENSO-induced supply disruption and hedge
forward accordingly, the volatility footprint of an ENSO anomaly should
appear near-contemporaneously with the climate signal, not lagged by the
agronomic timescale. The lag-decay pattern is therefore consistent with
market anticipation collapsing the physical lag, rather than with the
absence of an underlying mechanism.

What anticipation alone cannot explain, however, is the episodic and
non-persistent character of the cumulative gain. A stable anticipation
channel should deliver a steady forecast improvement across decades of ENSO
cycles, not gains that materialise in specific supply-stress episodes and
revert in the intervening years. The El~Ni\~no episode around the turn of
the millennium, which produced a temporary gain at the centre of the
distribution that dissipated over the following decade, provides a second
illustration of this pattern: even a
real ENSO--supply alignment does not create a structural signal that
survives into subsequent climate cycles. The
plausible bridge between these two observations is the rare-disaster reading
proposed by \textcite{bonatoNinoNinaForecastability2023}: ENSO acts as a
proxy for the level of uncertainty over future supply, and risk-averse
inventory holders respond to high ENSO uncertainty by building stocks, which
in turn raises the convenience yield and amplifies the variance of returns.
Under this reading, ENSO is predictive of upper-tail volatility only when
the underlying supply system is already close to a stress threshold; in
quiet periods, the same Ni\~no~3.4 readings carry no incremental information
beyond the HAR lags. The 2022--2024 episode supplies precisely that stress
threshold via the CSSVD outbreak and the resulting structural supply
shortfall, and the ENSO regressor identifies the El~Ni\~no transition that
coincides with it. \textcite{bouriNinoForecastabilityOilprice2021} find a
structurally similar pattern for oil, where the predictive value of ENSO is
concentrated at specific horizons (two to four years) and emerges most
strongly under longer estimation windows that capture full climate cycles;
the horizon over which the signal becomes visible differs across markets,
but the regime-conditional character of the effect is shared.

This reading is also consistent with the cross-commodity null.
\textcite{ubilavaRoleNinoSouthern2018} reports that ENSO produces in-sample
evidence of predictability for selected non-agricultural commodities but
that the out-of-sample relationship typically fails to replicate, and that
the effect is most amplified for cocoa during El~Ni\~no episodes
specifically. The HAR-X-QR results extend this pattern: coffee and sugar
share the same equatorial belt and the same Pacific climate exposure, yet
the same regressor that produces cocoa's crisis-concentrated twelve-month
gain \emph{worsens} long-horizon coffee forecasts and produces only an
economically negligible gain on sugar. A structural climate
channel operating uniformly across tropical soft commodities is hard to
reconcile with this asymmetry; a regime-state indicator that happens to
align with one severe cocoa-specific episode is not.
\textcite{huangHighfrequencyApproachVaR2022} provide further support for
this reading by showing that HAR coefficients differ substantially across
market regimes; the ENSO coefficient, on this interpretation, is identified
predominantly in the regime where its sign and magnitude happen to align
with the realised volatility surprise.

The data therefore support a clear answer to RQ2: the ENSO index does not
improve cocoa volatility forecasts in any robust sense. The full-sample
twelve-month upper-tail gain of $3.3\,\%$ is the residual of a single
crisis episode superimposed on decades of mild deterioration, and in the
pre-crisis subsample the augmented model is strictly worse. A coherent
mechanism, combining market anticipation with a rare-disaster
amplification through inventory dynamics, can rationalise both why
whatever predictive content ENSO carries appears near-contemporaneously
and why it surfaces only under supply stress. But a regressor whose
contribution is confined to one episode, identifiable only ex post, has
no operational forecasting value. The honest reading is that HAR-X-QR
captures a sample-specific covariation between the ENSO state and the
cocoa supply crisis, rather than a replicable structural climate
channel.

%------------------------------------------------------------------------------
\subsection{Forecast Instability and the Crisis Episode}
\label{sec:disc-instability}
%------------------------------------------------------------------------------

The pattern just described is not an isolated peculiarity of cocoa. It is a
textbook instance of the forecast-instability phenomenon documented in the
financial-econometrics literature. \textcite{payeInstabilityReturnPrediction2006}
show that the discrepancy between strong in-sample predictability and weak
out-of-sample performance, repeatedly observed across return-prediction
studies, is best explained by structural instability in the underlying
predictive relationship: predictability is real but concentrated in
identifiable subperiods rather than uniform over time.
\textcite{rossiForecastingPresenceInstabilities2020} argues that instability
is the rule in macroeconomic and financial forecasting, not the exception,
and that diagnostic procedures aimed at locating the periods in which
predictability is concentrated should now be considered part of standard
practice.

The expanding-window diagnostic in Section~\ref{sec:res-robustness} is
designed to operate within exactly that tradition. It imposes no break date,
partitions no sample ex post, and traces the cumulative loss difference
recursively from the first feasible forecast onward. The diagnostic
identifies two episodes rather than one. The 1997--98 El~Ni\~no --- the
strongest on record at that time --- produced a temporary gain at the
median and $\tau = 0.75$ for long horizons that subsequently dissipated
over the 2000s, reaching near zero well before the 2022--2024 crisis and
failing to reappear in the intervening ENSO cycles. That dissipation is
itself evidence of regime-conditional behaviour: ENSO produced a real but
transient predictive gain when the cocoa supply system was under stress,
and carried no incremental information in the years that followed. The
finding that the cumulative ENSO gain at the twelve-month upper tail is
negative --- the augmentation harmful on net --- through the three decades
preceding 2022, and that the entire full-sample improvement accumulates
within the final two years, is diagnostically clean: both partitions
emerge from the data, not from the researcher. \textcite{poonForecastingVolatilityFinancial} warn that
in-sample parameter stability is no guarantee of out-of-sample stability,
and that time variation in parameter estimates is the critical issue in
forecast evaluation; the expanding-window construction addresses this
warning directly.

The Model Confidence Set complements rather than contradicts this picture.
\textcite{hansen2011model} stress that the MCS acknowledges the limitations
of the data: informative samples yield small confidence sets, uninformative
samples include many models. The wide long-horizon confidence sets --- all
models except historical volatility share $\mathcal{M}^*$ from three months
onward --- are therefore a statement about the limits of the data: averaged
over the full sample, the loss differentials at long horizons are small
relative to their sampling variation, and the MCS cannot, by construction,
separate a stable channel from a single episode that dominates the average
loss. Read together, the MCS result and the expanding-window diagnostic are
entirely consistent: the former says the full-sample averages barely
discriminate among models at long horizons; the latter says that what
little discrimination exists is shaped by one specific subperiod.

%------------------------------------------------------------------------------
\subsection{Coverage, Calibration, and the Discrimination Tradeoff}
\label{sec:disc-coverage}
%------------------------------------------------------------------------------

Two coverage results in Section~\ref{sec:res-har-qr} appear, on a first
reading, to be in tension with the quantile-loss rankings. The HAR-QR
violation rate at the twelve-month upper quantile rises to $14.63\,\%$
against a nominal $5\,\%$, and the HAR-X-QR violation rate rises further to
$15.89\,\%$, yet at that horizon HAR-X-QR achieves a modestly lower quantile
loss than HAR-QR. \textcite{christoffersen1998evaluating} provides the relevant
distinction: unconditional coverage and conditional coverage are separate
properties, and a forecast can satisfy the former while violating the
latter through clustered exceedances. On the non-overlapping subsequences
used for testing, the $LR_{uc}$ test rejects unconditional coverage at six
months for both models, and for HAR-X-QR additionally at one and three
months (Table~\ref{tab:christoffersen}); at twelve months the subsequence
is too short for a powerful test, so the 14--16\,\% empirical violation
rates are the more informative summary. The concentration of violations in
the 2022--2024 episode is directly visible in
Figure~\ref{fig:upper-tail} (panel b2), where realized volatility remains
above the $\tau = 0.95$ boundary for an extended stretch before the
forecast distribution adjusts, so the test statistics understate the true
calibration problem.

The apparent contradiction between higher violations and lower quantile loss
is resolved by recognising that the two metrics measure orthogonal
properties of a distributional forecast. \textcite{dimitriadisStatisticalInferenceScore2026}
formalise this distinction in a score-decomposition framework: average loss
under a strictly consistent scoring function rewards
\emph{discrimination}, the ability to distinguish high-risk from low-risk
states, whereas backtest hit rates assess \emph{calibration}, the ability
to match nominal coverage. The two are not jointly maximised. The ENSO
coefficient in HAR-X-QR raises the upper-quantile forecast during El~Ni\~no
episodes, which reduces the magnitude of large exceedances when realised
volatility spikes, and lowers it during La~Ni\~na, which produces more
small exceedances in otherwise calm periods. The asymmetric check function
penalises a small number of large under-predictions more heavily than a
larger number of small ones, so a model that re-allocates its boundary
toward the conditions that produce the most damaging exceedances can win on
quantile loss while losing on average violation rate. At the twelve-month
horizon this is exactly what happens: HAR-X-QR raises the boundary during
the crisis enough to reduce the largest exceedances, gaining $3.3\,\%$ in
quantile loss at the cost of a higher violation rate. The qualification
matters, however: at the one- to six-month horizons ENSO augmentation
worsens loss and coverage \emph{simultaneously}, so the
discrimination--calibration tradeoff rationalises the twelve-month result
only, not the augmentation in general.

\textcite{pattonVolatilityForecastComparison} provides a complementary
argument for prioritising the score over the hit rate in this setting.
Among the loss functions commonly used in volatility evaluation, only MSE
and QLIKE are robust to noise in the volatility proxy; coverage diagnostics
compare a binary indicator against the noisy proxy directly, and their
informativeness deteriorates as proxy noise rises with the forecast
horizon. The Yang-Zhang proxy becomes noisier as the averaging horizon
extends; the quantile-loss ranking is, on Patton's argument, more
trustworthy at twelve months than the violation-rate comparison. Which
property a user should prioritise depends on the use case: distributional
ranking favours the score, nominal-coverage compliance favours the
violation rate, and the two should be reported together.

%------------------------------------------------------------------------------
\subsection{Limitations and Future Research}
\label{sec:disc-limitations}
%------------------------------------------------------------------------------

Six limitations of the present specification deserve explicit mention.
First, the HAR-QR coefficients are static within each estimation window.
\textcite{gongRealizedVolatilityForecasting2017} show that a HAR
specification with time-varying sparsity, using normal-gamma autoregressive
priors, substantially outperforms standard HAR for agricultural commodity
futures; a time-varying HAR-QR with smoothly evolving quantile coefficients
is a natural next step and would be expected to reduce the coverage
degradation observed at long horizons. Second, the model has no explicit
regime structure. \textcite{luoForecastingRealizedVolatility2022}
demonstrate that infinite-Hidden-Markov regime switching delivers
substantial gains over benchmark HAR for agricultural commodities precisely
when macroeconomic uncertainty is high; the 2022--2024 episode is exactly
the kind of high-uncertainty regime that an IHM-HAR-QR would be designed to
model directly. Third, the model conditions on volatility lags but not on
realised jumps, skewness, or higher moments;
\textcite{andersen2007roughing} and \textcite{bonatoForecastingRealizedVolatility2024}
both show that these components carry forecast-relevant information beyond
what HAR lags capture, and the upper-tail coverage failures at long
horizons are a plausible signature of unmodelled jumps in the crisis
period.

Fourth, the climate exposure is summarised by a single index. The Indian
Ocean Dipole and the Atlantic Meridional Mode also influence West African
rainfall patterns and could disentangle Pacific from Atlantic
teleconnections in a multi-index extension. Fifth, the analysis uses ICE
London front-month rolling futures denominated in GBP only; ICE New York
contracts, basis dynamics between the two markets, and contract-structure
effects are not modelled, and a comparative two-exchange analysis would
provide a natural robustness extension. Sixth, the estimation window length
is fixed at 2{,}520 trading days. \textcite{bouriNinoForecastabilityOilprice2021}
report that the predictive value of ENSO on oil volatility emerges most
strongly under estimation windows of four to six years; the ten-year
window used here may dilute regime-specific coefficients and dampen
climate-driven episodes. An adaptive-window design that lets the data
choose the relevant historical horizon for each forecast is a methodological
extension worth pursuing.

Beyond these direct extensions, the broader research programme suggested by
the cocoa results points in three directions. A jump-augmented HAR-J-QR
specification would address the upper-tail coverage problem at the source.
A regime-switching IHM-HAR-QR would model directly what the present
specification absorbs through a single ENSO coefficient. A multi-predictor
extension combining the Ni\~no~3.4 index with realised inventory data,
shipping rates, or other proxies of physical supply tightness would test
whether the ENSO gain identified here can be disentangled from the broader
regime indicators that coincide with it. Each of these directions would
require both the data infrastructure and the methodological development
that lie outside the scope of a single master's thesis; together they
define the natural research agenda that follows from the findings reported
in Section~\ref{sec:results}.

The overall picture is consistent. HAR-based quantile regression provides a
statistically significant improvement over the standard
volatility-forecasting benchmarks for cocoa at the daily-to-monthly
horizons, where it holds the top rank in an upper-tail Model Confidence
Set from which all three benchmarks are excluded. At the six- and
twelve-month horizons that motivated the
procurement perspective, simpler quantile constructions anchored to the
unconditional distribution prove more robust, and no model achieves
acceptable upper-tail calibration through the crisis. The ENSO augmentation
produces no robust gain at any horizon: its modest full-sample twelve-month
improvement reflects a single severe episode rather than a replicable
climate channel. The combination of a clearly delimited RQ1 result and a
negative RQ2 result positions the thesis as a methodological contribution
to the volatility-forecasting literature on agricultural commodities, while
flagging precisely where the limits of what daily-frequency data and three
decades of evaluation can reveal about a slow-moving climate signal
in fact lie.
